\begin{enquestion}{20}{
    The Royal Netherlands Army Command (CLAS) is testing a new type of battery for portable communication devices.
    The manufacturer claims an average service life of $\mu_0 = 36$ hours per fully charged battery.
    From previous production runs, it is known that the service life is normally distributed.
    A test batch of $n = 25$ batteries yields a mean of $35.6$ hours and a standard deviation of $3.04$ hours.
}

    \subquestion{5}{
        Calculate a two-sided \SI{96}{\percent}-confidence interval for the mean service life $\mu$ based on the given sample data.
        Round the boundaries to one decimal place in such a way that the confidence level is guaranteed.
        Is there a reason to believe that the mean service life differs significantly from $36$ hours?
    }

    \ensolution{
        The problem states that $\mu_0 = 36$ hours, the sample mean
        $\samplemeanx = 35.6$ hours, the sample standard deviation
        $s = 3.04$ hours, and the sample size $n = 25$.

        The service life $X$ of an arbitrary battery is normally distributed with $\mu = \text{ ?}$ and $\sigma = \text{ ?}$. 
        We want to calculate a confidence interval for the unknown $\mu$ based on the sample data. 
        
        Since the standard deviation $\sigma$ is unknown, we need to resort to the $t$-distribution for computing a confidence interval and use the sample standard deviation as an estimate.\rubric{1}
        
        Furthermore, $\alpha = 0.04$, because we want to calculate a \SI{96}{\percent} confidence interval.
        The corresponding $t$-value is equal to
        \[
            t_{n-1, \alpha/2} = \invnorm(\opp=1-\alpha/2=0.98, \df=n-1=24) \approx 2.1715.\rubric{1}
        \]
        The \SI{96}{\percent} confidence interval is then given by
        \begin{align*}
            &\phantom{=}[\samplemeanx - t_{n-1, \alpha/2} \cdot \frac{ s }{ \sqrt{ n } }, \samplemeanx + t_{n-1, \alpha/2} \cdot \frac{ s }{ \sqrt{ n } }] \\
            &= [35.6 - 2.1715 \cdot \frac{ 3.04 }{ \sqrt{ 25 } }, 35.6 + 2.1715 \cdot \frac{ 3.04 }{ \sqrt{ 25 } }] \\
            &= [34.2797, 36.9203].\rubric{2}
        \end{align*}
        Since the original value $\mu_0 = 36$ hours lies within the confidence interval, there is insufficient evidence to say that the mean service life differs significantly from $36$ hours.\rubric{1}
    }

    \subquestion{7}{
        The manufacturer wants to estimate the mean service life more precisely, namely with a maximum deviation of $\pm 0.5$ hours at a 
        \SI{96}{\percent} confidence level.
        What is the minimum required sample size $n$ to estimate with this precision?
    }

    \ensolution{
        Our first step is to obtain an initial value for $n$ by pretending that $\sigma$ is equal to the estimate $s$.
        If $\sigma$ would be known, we could use the corresponding $z$-value:
        \[
            z_{\alpha/2} = \invnorm(\opp=1-\alpha/2, \mu=0, \sigma=1) \approx 2.0537.\rubric{1}
        \]
        In that case, we can use the following formula:
        \begin{align*}
            n   &\geq \left( \frac{z_{\alpha/2} \cdot \sigma}{E} \right)^2. \\
                &= \left( \frac{2.0537 \cdot 3.04}{0.5} \right)^2 \\
                &\approx 155.9126. \rubric{2}
        \end{align*}
        We can therefore use $n = 156$ as an initial value for our search.

        Now, make a table on the graphical calculator of the expression:
        \[
            y_1 = \invt(\opp=1-\alpha/2=0.98, \df=X-1) \cdot 3.04 / \sqrt{X},
        \]
        and set the table range from $n=156$ to $n=165$ (most of the time the outcomes will only be slightly larger than the initial value).\rubric{1}
        We get:
        \begin{align*}
            X = 158: &\quad \frac{3.04}{\sqrt{158}} \cdot \invt(\opp=0.98, \df=158-1) \approx 0.5009. \\
            X = 159: &\quad \frac{3.04}{\sqrt{159}} \cdot \invt(\opp=0.98, \df=159-1) \approx 0.4993. \rubric{2}
        \end{align*}
        With other words, we need a sample with at least $n = 159$ batteries to be able to construct a \SI{96}{\percent}-confidence interval with error margin at most $\pm 0.5$ hours.\rubric{1}
    }

    \subquestion{8}{
        CLAS is not satisfied with the quality of the batteries from this supplier, and therefore places an order with a second supplier.
        For a test, a batch of $n = 64$ batteries are ordered from this new supplier.
        From this sample, it is determined that the mean service life is equal to $36.5$ hours and a standard deviation of $3.67$ hours.
        
        CLAS is primarily interested in whether the batteries last \emph{longer} than $36$ hours on average.
        Perform an appropriate hypothesis test and determine the test result using the critical region.
        Use a significance level of $\alpha = 0.04$.        
    }

    \ensolution{
        \subsubsection*{Step 1: formulate the null hypothesis $H_0$ and the alternative hypothesis $H_1$.}
        For a test of whether the mean $\mu$ of a normal distribution is larger than $\mu_0=36$, we use the following hypotheses
        \begin{align*}
            H_0: \quad \mu \leq 36\\
            H_1: \quad \mu > 36 \rubric{1}
        \end{align*}

        \subsubsection*{Step 2: determine the significance level $\alpha$.}
        The problem states that we may work with a significance level of $\alpha=0.04$.
        
        \subsubsection*{Step 3: gather data for calculating the test statistic.}
        Furthermore, sample data has already been gathered in the problem for the population of batteries from the second supplier, namely
        $\overline{x} = 36.5$, $s = 3.67$ hours and $n = 64$. 

        \subsubsection*{Step 4: calculate the test statistic.}
        The theoretical test statistic (capital letter!) for a $t$-test ($\sigma$ is unknown!) is 
        \[
            T = \frac{\overline{X}-\mu_0}{S / \sqrt{n}}. \rubric{1}
        \]
    
        Under $H_0$, this test statistic follows a $t$-distribution with $\df = n - 1 = 63$ degrees of freedom. \rubric{1}
        
        The observed test statistic (lowercase letter!) is equal to
        \begin{align*}
            t = \frac{\overline{x} - \mu_0}{s / \sqrt{n}} = \frac{36.5 - 36}{3.67 / \sqrt{64}} \approx 1.0899.\rubric{1}
        \end{align*}    
        
        \subsubsection*{Step 5: determine the test result (using the critical region).}

        Since we are performing a right-sided test, the critical region is of the form $(g, \infty)$, where
        \begin{align*}
            g = \invt(\opp=1-\alpha=0.96, \df=63) \approx 1.7794.
        \end{align*}       
        Thus, the critical region is equal to $(1.7794, \infty)$. \rubric{2}

        The observed test statistic $x = 36.5$ does not lie in the critical region,
        so we decide not to reject $H_0$.\rubric{1}

        \subsubsection*{Step 6: translate the test result back to the original context of the problem.}
        There is insufficient reason based on this sample to assume that the mean service life $\mu$
        of batteries from the second supplier is significantly longer than $36$ hours.\rubric{1}
    }
\end{enquestion}
